\documentclass[a4paper,12pt]{article}
\usepackage[utf8]{inputenc}
\usepackage[russian]{babel}
\usepackage{csquotes}
\usepackage{amsmath, amssymb, amsfonts}
\usepackage{geometry}
\geometry{top=2cm, bottom=2cm, left=2cm, right=2cm}
\usepackage{enumitem}
\setlength{\parindent}{0pt}
\setlength{\parskip}{0.8ex}

\title{%
  Теоретический минимум 2: Возвращение Ларисы \\
  Линейная алгебра наносит ответный удар \\[2ex]
  \large \enquote{Дифуры решать, мать уважать, \\ линал пересдавать, отца не забывать!}%
}
\author{
}
\date{}

\newcommand{\R}{\mathbb{R}}
\newcommand{\C}{\mathbb{C}}
\newcommand{\la}{\langle}
\newcommand{\ra}{\rangle}
\newcommand{\proj}{\operatorname{proj}}
\newcommand{\rg}{\operatorname{rg}}
\newcommand{\tr}{\operatorname{Tr}}
\newcommand{\im}{\operatorname{Im}}
\newcommand{\ker}{\operatorname{Ker}}

\begin{document}
\maketitle
\newpage

\section*{Ответы на вопросы (1--121)}


\subsection*{1. !!! Сформулируйте аксиому положительной определенности скалярного произведения в вещественном евклидовом пространстве.}
Для любого вектора \(x\) выполняется \(\langle x, x \rangle \ge 0\), причём \(\langle x, x \rangle = 0 \iff x = 0\).

\subsection*{2. !!! Сформулируйте неравенство Коши—Буняковского—Шварца для произвольного евклидова пространства.}
\[ |\langle x, y \rangle| \le \sqrt{\langle x, x \rangle} \cdot \sqrt{\langle y, y \rangle}. \]
Равенство \(\iff\) векторы \(x\) и \(y\) линейно зависимы.

\subsection*{3. !!! Дайте определение эрмитовой симметрии скалярного произведения в комплексном унитарном пространстве.}
\[ \langle x, y \rangle = \overline{\langle y, x \rangle} \quad \forall x, y. \]

\subsection*{4. !!! Сформулируйте аксиому линейности скалярного произведения по первому аргументу.}
\[ \langle \alpha x + \beta y, z \rangle = \alpha \langle x, z \rangle + \beta \langle y, z \rangle \quad \forall \alpha, \beta \in \mathbb{R} \ (\text{или } \mathbb{C}). \]

\subsection*{5. !!! Сформулируйте свойство полулинейности скалярного произведения по второму аргументу в комплексном евклидовом пространстве.}
\[ \langle x, \alpha y + \beta z \rangle = \overline{\alpha} \langle x, y \rangle + \overline{\beta} \langle x, z \rangle. \]

\subsection*{6. !!! Сформулируйте аксиому однородности для нормы в линейном пространстве.}
\[ \|\alpha x\| = |\alpha| \cdot \|x\| \quad \forall \alpha \in \mathbb{R} \ (\text{или } \mathbb{C}), \ \forall x. \]

\subsection*{7. !!! Запишите формулу, выражающую евклидову норму векторов через скалярное произведение.}
\[ \|x\| = \sqrt{\langle x, x \rangle}. \]

\subsection*{8. Сформулируйте геометрический смысл неравенства треугольника для норм векторов.}
Длина любой стороны треугольника не превосходит суммы длин двух других сторон: \(\|x+y\| \le \|x\| + \|y\|\).

\subsection*{9. !!! Сформулируйте аксиому положительной определенности нормы.}
\(\|x\| \ge 0\), и \(\|x\| = 0 \iff x = 0\).

\subsection*{10. Запишите тождество параллелограмма в терминах норм.}
\[ \|x+y\|^2 + \|x-y\|^2 = 2(\|x\|^2 + \|y\|^2). \]

\subsection*{11. !!! Сформулируйте неравенство треугольника для евклидовой нормы в вещественном пространстве.}
\[ \|x+y\| \le \|x\| + \|y\|. \]

\subsection*{12. Сформулируйте тождество параллелограмма Иордана — фон Неймана в нормированном пространстве.}
Нормированное пространство является евклидовым (норма порождается скалярным произведением) \(\iff\) выполняется тождество параллелограмма.

\subsection*{13. В чем заключается геометрический смысл того, что нормированное пространство является евклидовым?}
В нём определены углы между векторами: \(\cos \angle(x,y) = \frac{\langle x,y \rangle}{\|x\|\|y\|}\).

\subsection*{14. Верно ли, что любое подпространство евклидова пространства само является евклидовым с тем же скалярным произведением?}
Да, ограничение скалярного произведения на подпространство сохраняет все аксиомы.

\subsection*{15. !!! Дайте определение ортогональности двух векторов x и y в вещественном евклидовом пространстве.}
\(\langle x, y \rangle = 0\).

\subsection*{16. Сформулируйте геометрическую формулировку теоремы Пифагора в терминах норм векторов.}
Если \(x \perp y\), то \(\|x+y\|^2 = \|x\|^2 + \|y\|^2\).

\subsection*{17. !!! Дайте определение ортогональной системы векторов. Какой дополнительный признак требуется, чтобы система называлась ортонормированной?}
Система \(\{e_i\}\) ортогональна, если \(\langle e_i, e_j \rangle = 0\) при \(i \ne j\). Ортонормирована, если дополнительно \(\|e_i\| = 1\) для всех \(i\).

\subsection*{18. Может ли ортогональная система векторов содержать нулевой вектор? Обоснуйте одним предложением.}
Да, формально ортогональность сохраняется, но обычно нулевой вектор исключают, так как он не участвует в формировании базиса.

\subsection*{19. !!! Сформулируйте теорему о линейной независимости ортогональной системы векторов.}
Любая система ненулевых попарно ортогональных векторов линейно независима.

\subsection*{20. !!! Запишите условие ортогональности двух функций f(x) и g(x) в пространстве C[a, b] со стандартным скалярным произведением.}
\[ \int_a^b f(x) g(x) \, dx = 0. \]

\subsection*{21. !!! Дайте определение ортонормированного базиса (ОНБ) в конечномерном евклидовом пространстве.}
Базис, состоящий из попарно ортогональных векторов единичной длины.

\subsection*{22. !!! Запишите формулу для вычисления i-й координаты вектора x в ОНБ {e_1, e_2, …, e_n}.}
\[ x_i = \langle x, e_i \rangle. \]

\subsection*{23. Сформулируйте геометрический смысл равенства Парсеваля.}
Квадрат длины вектора равен сумме квадратов его координат в ОНБ (обобщённая теорема Пифагора).

\subsection*{24. !!! Запишите равенство Парсеваля для вещественного евклидова пространства.}
\[ \|x\|^2 = \sum_{i=1}^n \langle x, e_i \rangle^2. \]

\subsection*{25. Как связаны понятия ортогональной системы векторов и ортогонального базиса?}
Ортогональный базис — это максимальная (по числу векторов) ортогональная система, которая является базисом.

\subsection*{26. !!! Запишите евклидову норму вектора через его координаты в ортонормированном базисе.}
\[ \|x\| = \sqrt{\sum_{i=1}^n x_i^2}. \]

\subsection*{27. !!! Дайте определение матрицы Грама для системы векторов {x₁, …, xₖ}.}
Матрица \(G_{ij} = \langle x_i, x_j \rangle\).

\subsection*{28. Сформулируйте свойство симметричности матрицы Грама для вещественного евклидова пространства.}
\(G^T = G\) (симметрична).

\subsection*{29. !!! Чему равен определитель Грама, если среди векторов системы есть хотя бы два коллинеарных?}
Нулю, так как система линейно зависима.

\subsection*{30. Запишите связь между площадью параллелограмма, натянутого на два вектора x и y, и их определителем Грама.}
\[ S^2 = \det G(x,y) = \|x\|^2 \|y\|^2 - \langle x, y \rangle^2. \]

\subsection*{31. !!! Каково взаимное расположение векторов, если их матрица Грама оказалась строго диагональной?}
Векторы попарно ортогональны.

\subsection*{32. Как изменится определитель Грама системы векторов, если один из векторов умножить на число α?}
Умножится на \(\alpha^2\).

\subsection*{33. !!! Запишите формулу для первого шага алгоритма Грама — Шмидта: как выражается первый вектор ортогональной системы g₁ через первый вектор исходной системы f₁?}
\(g_1 = f_1\) (затем нормируют: \(e_1 = g_1 / \|g_1\|\)).

\subsection*{34. !!! Сформулируйте теорему о существовании ортонормированного базиса в произвольном конечномерном евклидовом пространстве.}
Любое конечномерное евклидово пространство обладает ортонормированным базисом (процесс Грама–Шмидта).

\subsection*{35. В чем заключается геометрический смысл вычитания проекции в ходе ортогонализации?}
Вычитание проекции на уже построенные ортогональные векторы даёт компоненту, ортогональную их линейной оболочке.

\subsection*{36. Дайте определение ортогонального дополнения L⊥ к подпространству L в евклидовом пространстве E.}
\[ L^\perp = \{ x \in E \mid \langle x, y \rangle = 0 \ \forall y \in L \}. \]

\subsection*{37. Запишите определение ортогональной проекции вектора x на подпространство L.}
Вектор \( \proj_L x \in L\) такой, что \(x - \proj_L x \in L^\perp\).

\subsection*{38. Какова размерность ортогонального дополнения L⊥, если dim E = n, a dim L = k?}
\(\dim L^\perp = n - k\).

\subsection*{39. !!! Дайте определение расстояния d(x, L) от вектора x до подпространства L в нормированном пространстве.}
\[ d(x, L) = \inf_{y \in L} \|x - y\|. \]

\subsection*{40. Какой вектор из подпространства L реализует минимум расстояния от вектора x до этого подпространства?}
Ортогональная проекция \( \proj_L x \).

\subsection*{41. Сформулируйте геометрический смысл теоремы о кратчайшем расстоянии в евклидовом пространстве.}
Кратчайшее расстояние достигается на ортогональной проекции; вектор отклонения перпендикулярен подпространству.

\subsection*{42. !!! Чему равен образ нулевого вектора A(0_V) при действии линейного оператора?}
Нулевому вектору \(0_W\) (свойство линейности: \(A(0)=A(0\cdot v)=0\cdot A(v)=0\)).

\subsection*{43. Приведите один пример нелинейного отображения в пространстве R².}
\(f(x,y) = (x^2, y)\) или \(f(x,y) = (\sin x, y)\).

\subsection*{44. !!! Дайте определение ядра ker A линейного оператора A: V → W.}
\[ \ker A = \{ v \in V \mid A v = 0_W \}. \]

\subsection*{45. !!! Что называют рангом (rg A) линейного оператора? Запишите связь этого понятия с образом оператора.}
\(\rg A = \dim \im A\).

\subsection*{46. !!! Сформулируйте теорему о связи размерностей ядра, образа и исходного пространства V для линейного оператора.}
\[ \dim V = \dim \ker A + \dim \im A. \]

\subsection*{47. !!! Сформулируйте алгоритм заполнения j-го столбца матрицы линейного оператора: координаты какого вектора и в каком базисе должны там находиться?}
В j-м столбце стоят координаты образа j-го базисного вектора в базисе пространства-образа.

\subsection*{48. !!! Запишите матричное уравнение, связывающее столбец координат прообраза X и столбец координат образа Y через матрицу оператора A.}
\[ Y = A \cdot X. \]

\subsection*{49. Дайте определение ранга линейного оператора.}
Размерность образа (максимальное число линейно независимых векторов-образов).

\subsection*{50. Сформулируйте условие обратимости линейного оператора A: V → V через ранг его матрицы.}
\(\rg A = \dim V\) (матрица невырождена, \(\det A \ne 0\)).

\subsection*{51. !!! Чему равен определитель матрицы линейного оператора, если этот оператор обладает нетривиальным ядром (ker A != \{0_v\})?}
\(\det A = 0\).

\subsection*{52. !!! Дайте определение линейного эндоморфизма пространства V.}
Линейный оператор \(A: V \to V\) (отображение пространства в себя).

\subsection*{53. Запишите определение композиции (произведения) двух линейных операторов A и B.}
\((AB)(x) = A(B(x))\).

\subsection*{54. !!! Дайте определение невырожденного линейного оператора через его ядро.}
\(\ker A = \{0_v\}\) (оператор инъективен, для эндоморфизма это эквивалентно биективности).

\subsection*{55. Как связана матрица обратного оператора $A^{-1}$ с матрицей исходного оператора A в фиксированном базисе?}
Матрица обратного оператора равна обратной матрице: \([A^{-1}] = [A]^{-1}\).

\subsection*{56. Сформулируйте свойство инъективности линейного оператора.}
\(A(x_1) = A(x_2) \Rightarrow x_1 = x_2\) (эквивалентно \(\ker A = \{0\}\)).

\subsection*{57. Сформулируйте алгоритм заполнения i-го столбца матрицы перехода S от старого базиса к новому.}
В i-м столбце стоят координаты i-го вектора нового базиса в старом базисе.

\subsection*{58. !!! Запишите формулу связи матриц одного и того же линейного оператора в старом (A) и новом ($A^'$) базисах.}
\[ A' = S^{-1} A S, \] где \(S\) — матрица перехода от старого базиса к новому.

\subsection*{59. Перечислите три основных инварианта подобных матриц.}
След, определитель, характеристический многочлен (и его корни — собственные значения).

\subsection*{60. !!! Дайте определение собственного вектора линейного оператора. Могут ли входить в это множество нулевые векторы?}
\(v \ne 0\) такой, что \(A v = \lambda v\). Нулевой вектор не считается собственным.

\subsection*{61. !!! Запишите характеристическое уравнение, используемое для нахождения собственных значений матрицы линейного оператора.}
\[ \det(A - \lambda E) = 0. \]

\subsection*{62. Какова связь между следом матрицы оператора (Tr A) и корнями его характеристического многочлена в комплексной области?}
След равен сумме всех собственных значений (с учётом кратностей).

\subsection*{63. !!! Дайте определение собственного подпространства V_λ, отвечающего собственному значению λ.}
\[ V_\lambda = \{ v \in V \mid A v = \lambda v \}. \]

\subsection*{64. Что называют геометрической кратностью собственного значения λ? Свяжите это понятие с дефектом матрицы (A − λI).}
Геометрическая кратность = \(\dim V_\lambda = n - \rg(A - \lambda I)\) (дефект).

\subsection*{65. Сформулируйте необходимый и достаточный критерий диагонализуемости матрицы линейного оператора в терминах базиса.}
Существует базис из собственных векторов.

\subsection*{66. Дайте определение диагонализуемого линейного оператора.}
Оператор, матрица которого в некотором базисе диагональна.

\subsection*{67. Как связаны между собой элементы на главной диагонали диагональной матрицы оператора и его собственные значения?}
Они совпадают (с учётом кратностей).

\subsection*{68. !!! Запишите матричное равенство, выражающее исходную матрицу A через её диагональный вид D и матрицу перехода S, состоящую из собственных векторов.}
\[ A = S D S^{-1}. \]

\subsection*{69. !!! Запишите тождество, определяющее сопряженный оператор A* для линейного оператора A в евклидовом пространстве.}
\[ \langle A x, y \rangle = \langle x, A^* y \rangle \quad \forall x, y. \]

\subsection*{70. !!! Какова связь между матрицей оператора A и матрицей его сопряженного оператора A* в произвольном ортонормированном базисе (ОНБ)?}
\[ [A^*] = [A]^T \quad (\text{в вещ. случае}); \quad [A^*] = [A]^{\overline{T}} \ (\text{комплексный}). \]

\subsection*{71. Сформулируйте теорему Рисса о представлении линейного функционала в евклидовом пространстве.}
Любой линейный функционал \(f(x)\) имеет вид \(f(x) = \langle x, y_f \rangle\) для единственного \(y_f\).

\subsection*{72. !!! Дайте определение самосопряженного линейного оператора в вещественном евклидовом пространстве.}
\(A^* = A\) (т.е. \(\langle A x, y \rangle = \langle x, A y \rangle\)).

\subsection*{73. !!! Сформулируйте критерий самосопряженности оператора в терминах его матрицы в ортонормированном базисе (ОНБ).}
Матрица симметрична: \([A]^T = [A]\).

\subsection*{74. Чему равно скалярное произведение ⟨A��, ��⟩ для самосопряженного оператора A, если выразить его через действие оператора на второй аргумент?}
\(\langle A x, y \rangle = \langle x, A y \rangle\).

\subsection*{75. Какими числами (вещественными или комплексными) всегда являются собственные значения самосопряженного оператора?}
Только вещественными.

\subsection*{76. Сформулируйте геометрическое свойство взаимного расположения двух собственных векторов самосопряженного оператора, если они отвечают различным собственным числам.}
Они ортогональны.

\subsection*{77. Запишите комплексное эрмитово свойство скалярного произведения, используемое при доказательстве вещественности спектра: ⟨λx, x⟩ через вынесение скаляра.}
\(\langle \lambda x, x \rangle = \lambda \langle x, x \rangle\), а \(\langle x, \lambda x \rangle = \overline{\lambda} \langle x, x \rangle\). Из равенства (для самосопряжённого) следует \(\lambda = \overline{\lambda}\).

\subsection*{78. Сформулируйте спектральную теорему о существовании базиса из собственных векторов для самосопряженного оператора. Какое геометрическое свойство имеет этот базис?}
Существует ортонормированный базис из собственных векторов.

\subsection*{79. Запишите формулу спектрального разложения линейного самосопряженного оператора через сумму проекторов на его собственные подпространства.}
\[ A = \sum_i \lambda_i P_i, \quad P_i \text{ — ортопроектор на } V_{\lambda_i}. \]

\subsection*{80. Что представляет собой оператор проектора P = v vᵀ на направление нормированного собственного вектора v в координатном представлении?}
Матрица \(P_{ij} = v_i v_j\); \(P x = \langle x, v \rangle v\).

\subsection*{81. !!! Дайте определение ортогонального оператора в вещественном евклидовом пространстве.}
Оператор, сохраняющий скалярное произведение: \(\langle Ux, Uy \rangle = \langle x, y \rangle\).

\subsection*{82. !!! Какое матричное соотношение связывает ортогональную матрицу Q и её транспонированную матрицу Qᵀ в ортонормированном базисе?}
\[ Q^T Q = Q Q^T = E. \]

\subsection*{83. !!! Какие значения может принимать определитель матрицы ортогонального оператора?}
\(\det Q = \pm 1\).

\subsection*{84. Дайте определение билинейной формы, заданной на линейном пространстве V.}
Отображение \(B: V \times V \to \mathbb{R}\) (или \(\mathbb{C}\)), линейное по каждому аргументу.

\subsection*{85. Запишите формулу связи матриц одной и той же билинейной формы в старом (A) и новом ($A^'$) базисах при матрице перехода C.}
\[ A' = C^T A C. \]

\subsection*{86. Напишите формулу поляризации, выражающую симметричную билинейную форму B(x, y) через порожденную ею квадратичную форму Q(x).}
\[ B(x,y) = \frac{1}{2}\big(Q(x+y) - Q(x) - Q(y)\big), \quad Q(x)=B(x,x). \]

\subsection*{87. !!! Дайте определение канонического вида квадратичной формы.}
Диагональный вид, где коэффициенты при квадратах равны \(0, \pm 1\) (или любые ненулевые числа, в зависимости от соглашения).

\subsection*{88. Что называют сигнатурой и рангом квадратичной формы?}
Ранг — число ненулевых коэффициентов в каноническом виде. Сигнатура — пара \((\pi, \nu)\) (количество положительных и отрицательных квадратов).

\subsection*{89. !!! Сформулируйте закон инерции Сильвестра для квадратичных форм над вещественным полем.}
Число положительных и отрицательных квадратов в каноническом виде не зависит от способа приведения.

\subsection*{90. Дайте определение строго положительно определенной квадратичной формы.}
\(Q(x) > 0\) для всех \(x \ne 0\).

\subsection*{91. !!! Сформулируйте критерий Сильвестра для строго отрицательно определенной квадратичной формы (укажите знаки главных миноров).}
Квадратичная форма отрицательно определена \(\iff\) все главные миноры чередуют знак, начиная с отрицательного: \(\Delta_1 < 0,\ \Delta_2 > 0,\ \Delta_3 < 0, \dots\)

\subsection*{92. !!! Дайте определение особого решения обыкновенного дифференциального уравнения первого порядка.}
Особым решением называется такое решение, которое не может быть получено из общего решения ни при каких значениях произвольной постоянной (в том числе и при предельных переходах), но во всех его точках нарушается условие единственности решения задачи Коши.
\subsection*{93. !!! Сформулируйте задачу Коши для ОДУ первого порядка, разрешенного относительно производной.}
Найти решение \(y = y(x)\) уравнения \(y' = f(x,y)\), удовлетворяющее условию \(y(x_0) = y_0\).

\subsection*{94. !!! Запишите условие Липшица по переменной y для функции f(x, y) на некотором промежутке.}
\[ |f(x, y_1) - f(x, y_2)| \le L |y_1 - y_2| \quad \forall x \in [a,b],\ \forall y_1, y_2. \]

\subsection*{95. !!! Укажите стандартную замену переменной для понижения порядка уравнения вида y'' = f(x, y'). Как при этом выражается вторая производная y''?}
Замена \(p = y'\), тогда \(y'' = p'\). Получим \(p' = f(x, p)\).

\subsection*{96. !!! Укажите стандартную замену переменной для понижения порядка уравнения вида y'' = f(y, y'). Как при этом выражается вторая производная y'' по правилу дифференцирования сложной функции?}
Замена \(p = y' = p(y)\), тогда \(y'' = p \frac{dp}{dy}\). Получим \(p \frac{dp}{dy} = f(y, p)\).

\subsection*{97. !!! Сколько независимых произвольных постоянных $$(C_1, C_2, …)$содержит общее решение обыкновенного дифференциального уравнения второго порядка?}
Две (\(C_1, C_2\)).

\subsection*{98. !!! Дайте определение линейно независимой системы функций $\{y_1(x), y_2(x), …, y_n(x)\}$ на интервале (a, b).}
Линейная комбинация \(\sum c_i y_i(x) \equiv 0\) на \((a,b)\) возможна только при всех \(c_i = 0\).

\subsection*{99. !!! Запишите определитель Вронского (Вронскиан) для системы двух функций y₁(x) и y₂(x).}
\[ W(x) = \begin{vmatrix} y_1(x) & y_2(x) \\ y_1'(x) & y_2'(x) \end{vmatrix} = y_1 y_2' - y_2 y_1'. \]

\subsection*{100. !!! Сформулируйте определение фундаментальной системы решений (ФСР) для линейного однородного уравнения n-го порядка.}
Совокупность \(n\) линейно независимых решений.

\subsection*{101. Запишите формулу Остроградского—Лиувилля для линейного однородного дифференциального уравнения второго порядка y'' + p(x)y' + q(x)y = 0.}
\[ W(x) = W(x_0) \exp\!\left(-\int_{x_0}^x p(t)\,dt\right). \]

\subsection*{102. Какое фундаментальное свойство Вронскиана решений следует из знака экспоненты в формуле Лиувилля?}
Вронскиан либо тождественно равен нулю, либо нигде не обращается в ноль.

\subsection*{103. Напишите формулу, позволяющую восстановить общее решение ЛОУ второго порядка, если известно одно его ненулевое частное решение y₁(x) и Вронскиан W(x).}
\[ y_2(x) = y_1(x) \int \frac{W(x)}{y_1^2(x)}\,dx. \]

\subsection*{104. !!! Запишите характеристическое уравнение для линейного однородного дифференциального уравнения второго порядка ay'' + by' + cy = 0.}
\[ a\lambda^2 + b\lambda + c = 0. \]

\subsection*{105. !!! Каков вид общего решения ЛОУ второго порядка, если характеристическое уравнение имеет один двукратный вещественный корень λ?}
\[ y(x) = C_1 e^{\lambda x} + C_2 x e^{\lambda x}. \]

\subsection*{106. Опираясь на подстановку Эйлера, объясните, почему оператор дифференцирования D превращает ЛОУ в алгебраический многочлен.}
При \(y = e^{\lambda x}\) получаем \(y^{(k)} = \lambda^k e^{\lambda x}\), и уравнение \(a_n y^{(n)} + \dots + a_0 y = 0\) переходит в \( (a_n \lambda^n + \dots + a_0) e^{\lambda x} = 0\).

\subsection*{107. !!! Запишите формулу Эйлера, связывающую комплексную экспоненту e^{iβx} с тригонометрическими функциями.}
\[ e^{i\beta x} = \cos(\beta x) + i \sin(\beta x). \]

\subsection*{108. !!! Каков вид общего вещественного решения ЛОУ второго порядка, если корни характеристического уравнения чисто мнимые: $\lambda_1,\lambda_2$ = ±iβ?}
\[ y(x) = C_1 \cos(\beta x) + C_2 \sin(\beta x). \]

\subsection*{109. Назовите физическую причину, которая превращает затухающие колебания системы в вечный гармонический осциллятор.}
Отсутствие сил трения (диссипации) — коэффициент затухания равен нулю.

\subsection*{110. Дайте определение фазовой траектории нормальной системы дифференциальных уравнений.}
Кривая в фазовом пространстве, параметризованная временем, каждая точка которой — состояние системы \((y_1(t), \dots, y_n(t))\).

\subsection*{111. !!! Запишите характеристическое уравнение для линейной системы Y' = A Y с постоянной матрицей A размера 2×2.}
\[ \det(A - \lambda E) = 0, \quad \text{т.е. } \lambda^2 - (\tr A)\lambda + \det A = 0. \]

\subsection*{112. !!! Каков вид векторного фундаментального решения Y₁(t), отвечающего вещественному собственному значению λ₁ и собственному вектору v⁽¹⁾?}
\[ Y_1(t) = e^{\lambda_1 t} v^{(1)}. \]

\subsection*{113. Сформулируйте связь между следом матрицы системы Tr(A) и знаком вещественной части её комплексных собственных значений α = Re(λ).}
\(\tr A = \lambda_1 + \lambda_2 = 2\alpha\) (для комплексно-сопряжённых), поэтому знак \(\tr A\) совпадает со знаком \(\alpha\).

\subsection*{114. Какой тип фазового портрета соответствует чисто мнимым собственным значениям матрицы системы λ₁,₂ = ±iβ?}
Центр (эллиптические траектории, устойчиво, но не асимптотически).

\subsection*{115. Напишите формулу Эйлера для раскрытия комплексной векторной экспоненты v e^{(α + iβ)t}.}
\[ v e^{\alpha t} (\cos(\beta t) + i \sin(\beta t)) = e^{\alpha t} \big( \Re(v) \cos(\beta t) - \Im(v) \sin(\beta t) + i(\dots) \big). \]

\subsection*{116. Какое алгебраическое соотношение определяет присоединенный вектор h, отвечающий кратному собственному значению λ и собственному вектору v?}
\[ (A - \lambda E) h = v. \]

\subsection*{117. В каком случае матрица системы второго порядка с кратным корнем λ имеет геометрическую кратность, равную 2? Какой вид имеет эта матрица?}
Когда \(A = \lambda E\) (скалярная матрица). Тогда \(\dim V_\lambda = 2\).

\subsection*{118. Какой тип фазового портрета возникает, если кратный корень характеристического уравнения строго отрицателен (λ < 0), а геометрическая кратность равна 1?}
Вырожденный устойчивый узел.

\subsection*{119. Сформулируйте геометрический смысл асимптотической устойчивости положения равновесия по Ляпунову.}
Все фазовые траектории, начинающиеся достаточно близко к равновесию, стремятся к нему при \(t \to +\infty\).

\subsection*{120. !!! Каким должно быть соотношение знаков вещественных частей всех собственных значений матрицы системы, чтобы точка равновесия была асимптотически устойчивой?}
Вещественные части всех собственных значений должны быть строго отрицательны.

\subsection*{121. !!! Что можно сказать об устойчивости линейной системы, если среди собственных значений её матрицы обнаружилось хотя бы одно число с положительной вещественной частью (Re(\lambda) > 0)?}
Положение равновесия неустойчиво.

\end{document}